% Lösungen Einheit 2 von 4 – Rechte Winkel rückwärts (zusammenbau v0.1)
\einheitenkopf{Einheit 2 von 4 · Rechte Winkel rückwärts}

\erg{14}{a) null setzen – ein Wert ist gesucht, Gleichung lösen \quad b) $\overrightarrow{BO} \circ \overrightarrow{BA} = (-2 | -2 | -3) \circ (2 | -2 | a - 3) = -4 + 4 - 3a + 9 = 9 - 3a = 0$, also $a = 3$ \quad c) $\overrightarrow{Q_tO} \circ \overrightarrow{Q_tP} = (-2 | -t | -8) \circ (2 | -t | 4) = -4 + t^2 - 32 = t^2 - 36 = 0$, also $t = 6$ oder $t = -6$ \quad d) $Q(4 | 0 | q)$: $\overrightarrow{QR} \circ \overrightarrow{QF} = (-4 | 3 | 6 - q) \circ (6 | 0 | 8 - q) = q^2 -14q + 24 = 0$, also $q = 2$ oder $q = 12$; nur $q = 2$ liegt auf der Kante ($0 \le q \le 10$) \quad e) $\overrightarrow{AG} \circ \overrightarrow{CE} = (-5 | 12 | h) \circ (5 | -12 | h) = -25 - 144 + h^2 = 0$, also $h = 13$ ($h = -13$ entfällt); $V = 5 \cdot 12 \cdot 13 = 780$ \quad f) $C(x | y | 5)$: $\overrightarrow{AB} \circ \overrightarrow{AC} = 3(x - 1) + 2(y - 2) + 2 \cdot 5 = 3x + 2y +3 = 0$ und $2x + y = 0$; Lösung $x = 3$, $y = -6$, also $C(3 | -6 | 5)$ \quad g) $T = \lambda \cdot \overrightarrow{AC}$ mit $0 \le \lambda \le 1$: $\overrightarrow{BA} \circ \overrightarrow{BT} = 0$ liefert $\lambda = \frac{1}{3}$; damit $|AT| : |CT| = \lambda : (1 - \lambda) = 1 : 2$ \quad h) $\overrightarrow{AC} = (a | 0 | c)$ mit $a + 0 + c = 0$, etwa $(1 | 0 | -1)$; $|\overrightarrow{AB}| = \sqrt{18}$, $|(1 | 0 | -1)| = \sqrt{2}$, Faktor $3$: $C(3 | 0 | -3)$ (oder $C(-3 | 0 | 3)$)}
\erg{15}{a) Fehler: Nele setzt die Vektoren von $O$ aus an und prüft damit den Winkel bei $O$. Richtig: $\overrightarrow{BO} \circ \overrightarrow{BA} = 9 - 3a = 0$, also $a = 3$ \quad b) Der rechte Winkel heißt Skalarprodukt der Schenkelvektoren gleich null; das Skalarprodukt ist ein Term in der unbekannten Koordinate, und „gleich null“ macht daraus eine Gleichung – eine Bedingung für eine Unbekannte. \quad c) $(-9 | 12 | h) \circ (9 | -12 | h) = -81 - 144 + h^2 = 0$, also $h = 15$ dm; $V = 9 \cdot 12 \cdot 15 = 1\,620$ dm³}
